\Article{13}
Die Hauptaufgabe besteht nun in einer Transformation der Funktion \(T\), ähnlich derjenigen, die wir in Theoria Motus Corporum Coelestium, Art. 182, und ausführlicher in Disquisitione de elementis ellipticis Palladis dargestellt haben. Wir setzen nämlich~(17)
\[
\begin{aligned}\relax
[bb,1]&=[bb]-\frac{[ab]^2}{[aa]}\\[4pt]
[bc,1]&=[bc]-\frac{[ab][ac]}{[aa]}\\[4pt]
[bd,1]&=[bd]-\frac{[ab][ad]}{[aa]}\\
\text{etc.}&\\
[cc,2]&=[cc]-\frac{[ac]^2}{[aa]}-\frac{[bc,1]^2}{[bb,1]}\\[4pt]
[cd,2]&=[cd]-\frac{[ac][ad]}{[aa]}-\frac{[bc,1][bd,1]}{[bb,1]}\\
\text{etc.}&\\
[dd,3]&=[dd]-\frac{[ad]^2}{[aa]}-\frac{[bd,1]^2}{[bb,1]}-\frac{[cd,2]^2}{[cc,2]}
\end{aligned}
\]
\SourceLeaf{075}
etc. etc. Setzen wir dann\footnote{Im Vorangehenden konnten für die verschiedenen Systeme von Größen, die auf die ersten drei Bedingungsgleichungen zu beziehen sind, jeweils drei Buchstaben genügen; hier erschien es jedoch angebracht, einen vierten hinzuzufügen, damit die Regel des Algorithmus deutlicher hervortritt. Da auf die Buchstaben \(a,b,c;\ A,B,C;\ \mathfrak A,\mathfrak B,\mathfrak C\) in der natürlichen Reihenfolge von selbst \(d,D,\mathfrak D\) folgen, in der Reihe \(x,y,z\) aber das Alphabet erschöpft ist, haben wir \(w\) angefügt und ebenso an die Reihe \(\xi,\eta,\zeta\) den Buchstaben \(\varphi\).}
\[
\begin{aligned}\relax
[bb,1]y+[bc,1]z+[bd,1]w+\text{etc.}&=\eta'\\
[cc,2]z.+[cd,2]w+\text{etc.}&=\zeta''\\
[dd,3]w+\text{etc.}&=\varphi'''
\end{aligned}
\]
\begin{center}etc., so ist\end{center}
\[
T=\frac{\xi\xi}{[aa]}+\frac{\eta'\eta'}{[bb,1]}+\frac{\zeta''\zeta''}{[cc,2]}+\frac{\varphi'''\varphi'''}{[dd,3]}+\text{etc.}
\]
und die Größen \(\eta',\zeta'',\varphi'''\) etc. hängen von \(\xi,\eta,\zeta,\varphi\) etc. durch die folgenden Gleichungen ab:
\[
\begin{aligned}
\eta'&=\eta-\frac{[ab]}{[aa]}\xi\\[5pt]
\zeta''&=\zeta-\frac{[ac]}{[aa]}\xi-\frac{[bc,1]}{[bb,1]}\eta'\\[5pt]
\varphi'''&=\varphi-\frac{[ad]}{[aa]}\xi-\frac{[bd,1]}{[bb,1]}\eta'-\frac{[cd,2]}{[cc,2]}\zeta''\\
\text{etc.}&
\end{aligned}
\]

Alle für unseren Zweck erforderlichen Formeln lassen sich nun leicht daraus entnehmen. Zur Bestimmung der Korrelate \(A,B,C\) etc. setzen wir nämlich (18)
\[
\begin{aligned}
\mathfrak B'&=\mathfrak B-\frac{[ab]}{[aa]}\mathfrak A\\[5pt]
\mathfrak C''&=\mathfrak C-\frac{[ac]}{[aa]}\mathfrak A-\frac{[bc,1]}{[bb,1]}\mathfrak B'\\[5pt]
\mathfrak D'''&=\mathfrak D-\frac{[ad]}{[aa]}\mathfrak A-\frac{[bd,1]}{[bb,1]}\mathfrak B'-\frac{[cd,2]}{[cc,2]}\mathfrak C''
\end{aligned}
\]
etc.; dann werden \(A,B,C,D\) etc. durch die folgenden Formeln ermittelt, und zwar in umgekehrter Reihenfolge, beginnend mit der letzten,
\[
\left.\begin{aligned}\relax
[aa]A+[ab]B+[ac]C+[ad]D+\text{etc.}&=\mathfrak A\\
[bb,1]B+[bc,1]C+[bd,1]D+\text{etc.}&=\mathfrak B'\\
[cc,2]C+[cd,2]D+\text{etc.}&=\mathfrak C''\\
[dd,3]D+\text{etc.}&=\mathfrak D'''
\end{aligned}\right\}\ (19)
\]
\begin{center}etc.\end{center}
\SourceLeaf{076}
Für die Summe \(S\) haben wir hingegen die neue Formel (20)
\[
S=\frac{\mathfrak A\mathfrak A}{[aa]}+\frac{\mathfrak B'\mathfrak B'}{[bb,1]}+\frac{\mathfrak C''\mathfrak C''}{[cc,2]}+\frac{\mathfrak D'''\mathfrak D'''}{[dd,3]}\quad\text{etc.}
\]
Wird schließlich das Gewicht \(P\) gesucht, das der plausibelsten Bestimmung der durch die Funktion \(u\) ausgedrückten Größe zuzuweisen ist, so setzen wir (21)
\[
\begin{aligned}\relax
[bl,1]&=[bl]-\frac{[ab][al]}{[aa]}\\[5pt]
[cl,2]&=[cl]-\frac{[ac][al]}{[aa]}-\frac{[bc,1][bl,1]}{[bb,1]}\\[5pt]
[dl,3]&=[dl]-\frac{[ad][al]}{[aa]}-\frac{[bd,1][bl,1]}{[bb,1]}-\frac{[cd,2][cl,2]}{[cc,2]}
\end{aligned}
\]
etc.; danach ist (22)
\[
\frac1P=[ll]-\frac{[al]^2}{[aa]}-\frac{[bl,1]^2}{[bb,1]}-\frac{[cl,2]^2}{[cc,2]}-\frac{[dl,3]^2}{[dd,3]}-\text{etc.}
\]

Die Formeln (17)\(.\ .\ .\ .\)(22), deren Einfachheit nichts zu wünschen übrig zu lassen scheint, liefern eine in jeder Hinsicht vollständige Lösung unseres Problems.

\Article{14}
Nachdem wir die Hauptprobleme erledigt haben, werden wir noch einige Nebenfragen berühren, die diesen Gegenstand weiter erhellen.

Zunächst ist zu untersuchen, ob die Elimination. durch die \(x,y,z\) etc. aus \(\xi,\eta,\zeta\) etc. zu ermitteln sind, jemals unmöglich werden kann. Dies wird offenbar eintreten, wenn die Funktionen \(\xi,\eta,\zeta\) etc. voneinander nicht unabhängig sind. Nehmen wir also vorübergehend an, eine von ihnen werde bereits durch die übrigen bestimmt, sodass die identische Gleichung
\[
\alpha\xi+\beta\eta+\gamma\zeta+\text{etc.}=0
\]
gilt, wobei \(\alpha,\beta,\gamma\) etc. bestimmte Zahlen bezeichnen. Somit ist
\[
\begin{aligned}
\alpha[aa]+\beta[ab]+\gamma[ac]+\text{etc.}&=0\\
\alpha[ab]+\beta[bb]+\gamma[bc]+\text{etc.}&=0\\
\alpha[ac]+\beta[bc]+\gamma[cc]+\text{etc.}&=0
\end{aligned}
\]
\Needspace{23\baselineskip}
etc.; setzen wir daher
\SourceLeaf{077}
\[
\begin{aligned}
\alpha a+\beta b+\gamma c+\text{etc.}&=p\theta\\
\alpha a'+\beta b'+\gamma c'+\text{etc.}&=p'\theta'\\
\alpha a''+\beta b''+\gamma c''+\text{etc.}&=p''\theta''
\end{aligned}
\]
\Needspace{12\baselineskip}
etc., so folgt von selbst
\[
\begin{aligned}
a\theta+a'\theta'+a''\theta''+\text{etc.}&=0\\
b\theta+b'\theta'+b''\theta''+\text{etc.}&=0\\
c\theta+c'\theta'+c''\theta''+\text{etc.}&=0
\end{aligned}
\]
etc. sowie
\[
p\theta\theta+p'\theta'\theta'+p''\theta''\theta''+\text{etc.}=0
\]
Da alle \(p,p',p''\) etc. ihrer Natur nach positive Größen sind, kann diese Gleichung offenbar nur bestehen, wenn \(\theta=0,\ \theta'=0,\ \theta''=0\) etc. gilt.

Betrachten wir nun die Werte der totalen Differentiale \(\mathrm dX,\mathrm dY,\mathrm dZ\) etc., die denjenigen Werten der Größen \(v,v',v''\) etc. entsprechen, auf die sich die Beobachtungen beziehen. Diese Differentiale, nämlich
\[
\begin{aligned}
&a\mathrm dv+a'\mathrm dv'+a''\mathrm dv''+\text{etc.}\\
&b\mathrm dv+b'\mathrm dv'+b''\mathrm dv''+\text{etc.}\\
&c\mathrm dv+c'\mathrm dv'+c''\mathrm dv''+\text{etc.}
\end{aligned}
\]
etc., hängen nach dem eben erreichten Schluss so voneinander ab, dass sie, jeweils mit \(\alpha,\beta,\gamma\) etc. multipliziert, eine identisch verschwindende Summe ergeben; oder, was dasselbe ist, jedes von ihnen, dessen zugehöriger Multiplikator \(\alpha,\beta,\gamma\) etc. nicht verschwindet, verschwindet von selbst, sobald alle übrigen als verschwindend vorausgesetzt werden. Daher ist mindestens eine der Bedingungsgleichungen \(X=0,\ Y=0,\ Z=0\) etc. als \textit{überflüssig} anzusehen, da sie von selbst erfüllt ist, sobald die übrigen erfüllt sind.

Betrachtet man die Sache jedoch genauer, so zeigt sich, dass dieser Schluss für sich genommen nur für einen unendlich kleinen Veränderungsbereich der Unbestimmten gilt. Genau genommen sind nämlich zwei Fälle zu unterscheiden: der eine, in dem eine der Bedingungsgleichungen \(X=0,\ Y=0,\ Z=0\) etc. schlechthin und allgemein bereits in den übrigen enthalten ist, was sich in jedem Fall leicht vermeiden lässt; der andere, in dem gleichsam zufällig bei denjenigen konkreten Werten der Größen \(v,v',v''\) etc., auf die sich die Beobachtungen bezie\SourceLeaf{078}%
hen, eine der Funktionen \(X,Y,Z\) etc., etwa die erste, \(X\), einen größten oder kleinsten Wert (oder allgemeiner einen stationären Wert) hinsichtlich aller Änderungen annimmt, die wir an den Größen \(v,v',v''\) etc. unter Wahrung der Gleichungen \(Y=0,\ Z=0\) etc. vornehmen könnten. Da jedoch in unserer Untersuchung die Veränderlichkeit der Größen nur innerhalb so enger Grenzen betrachtet wird, dass sie wie eine unendlich kleine behandelt werden kann, hat dieser zweite Fall (der in der Praxis kaum jemals vorkommen wird) dieselbe Wirkung wie der erste: Eine der Bedingungsgleichungen ist nämlich als überflüssig zu verwerfen. Wir können sicher sein, dass die Elimination notwendig möglich ist, wenn alle beibehaltenen Bedingungsgleichungen in dem hier gemeinten Sinn voneinander unabhängig sind. Eine ausführlichere Untersuchung müssen wir jedoch einer anderen Gelegenheit vorbehalten; dieser Gegenstand ist ihrer wegen seiner theoretischen Feinheit eher als wegen seines praktischen Nutzens nicht unwürdig.

\clearpage
